% =====================================================================
%  TITRE V : DE LA DISCIPLINE ET DU CONTRÔLE INTERNE
% =====================================================================
\titrepartie{Titre V}{De la discipline et du contrôle interne}

% ---------------------------------------------------------------------
\article{Application du règlement et régime des sanctions}

\cl[application]{L'ensemble des agents veille au respect du présent règlement. Les agents titulaires au minimum du grade de Senior Lead Officer peuvent prononcer des sanctions disciplinaires, dans la limite des sanctions ouvertes à leur échelon et en lien avec l'équipe de supervision.}

\cl[legalite]{Aucun agent ne peut être sanctionné d'une manière qui ne soit pas expressément prévue au présent article, ni pour un fait qui ne constitue pas la violation d'une disposition du présent règlement.}

\cl{Lorsqu'un fait fautif n'est visé par aucune disposition, une demande de sanction est adressée au Command Staff et à l'État-Major, qui statuent par note de service motivée précisant le fondement retenu et les éléments justificatifs.}

\sstitre{Échelle des sanctions et autorité compétente}

\vspace{2pt}
\noindent
\begin{tabular}{P{42mm} P{106mm}}
\headrow
\thead{Autorité} & \thead{Sanctions ouvertes} \\
\tkey{Senior Lead Officer} & \tcell{Avertissement.\newline Blâme, sous l'accord d'un membre de la supervision au minimum Sergeant I.} \\
\altrow
\tkey{Sergeant I et II} & \tcell{Avertissement.\newline Blâme de niveau I ou II.\newline Mise à pied de 0 à 72 heures.\newline Révocation, sous l'accord d'un membre du Command Staff.} \\
\tkey{Command Staff\newline et État-Major} & \tcell{Avertissement.\newline Blâme de niveau I ou II.\newline Mise à pied de durée déterminée ou indéterminée.\newline Rétrogradation.\newline Révocation.} \\
\bottomrule
\end{tabular}

\cl[cumul-sanctions]{Les sanctions se cumulent selon l'échelle suivante :}

\vspace{4pt}
\noindent
\begin{tabular}{C{36mm} C{16mm} P{92mm}}
\headrow
\thead{Cumul} & \thead{} & \thead{Sanction résultante} \\
\tcell{2 avertissements} & \tcell{\textbf{=}} & \tcell{1 blâme} \\
\altrow
\tcell{3 blâmes} & \tcell{\textbf{=}} & \tcell{1 révocation} \\
\bottomrule
\end{tabular}

\cl[prescription]{Lorsqu'aucune sanction n'est prononcée à l'encontre d'un agent pendant une période de trois mois, son casier disciplinaire est effacé.}

\cl[hierarchie-sanction]{À partir du grade de Lieutenant, le supérieur hiérarchique direct peut sanctionner son subordonné direct. Ainsi, le Captain sanctionne le Lieutenant, le Lieutenant sanctionne le Sergeant II.}

\cl[exclusivite]{La rétrogradation, la révocation et la promotion ne peuvent être prononcées que par le Command Staff et l'État-Major. Un superviseur peut en formuler la demande motivée.}

\cl[convocation]{L'agent convoqué à un entretien disciplinaire conserve ses fonctions jusqu'à la date de l'entretien. Une mise à pied conservatoire ne s'applique que si la convocation en fait expressément mention. En l'absence de mention, aucune mise à pied n'est prononcée. Lorsqu'elle est prononcée, cette mesure est conservatoire et ne préjuge d'aucune faute.}

\cl{Toute sanction est notifiée par écrit, motivée, et versée au dossier de l'agent. Une sanction non notifiée est réputée inexistante.}

\begin{note}[Proportionnalité de la sanction]
L'autorité qui sanctionne tient compte de la gravité objective du manquement, de ses conséquences, de l'intention de son auteur, de son ancienneté, de son casier disciplinaire et de son attitude après les faits. Une même faute ne donne pas nécessairement lieu à la même sanction pour tous.
\end{note}

% ---------------------------------------------------------------------
\article{Division des Affaires Internes et plaintes}

\chapeau{Un service de police qui ne se contrôle pas lui-même est contrôlé de l'extérieur, et dans de moins bonnes conditions. La Division des Affaires Internes existe autant pour sanctionner les fautes que pour établir l'innocence des agents mis en cause à tort.}

\cl[iab]{La Division des Affaires Internes est rattachée directement au Chief et demeure indépendante de la chaîne hiérarchique opérationnelle. Aucun de ses enquêteurs ne peut instruire une affaire concernant sa propre division ou une personne avec laquelle il entretient un lien personnel.}

\sstitre{Saisine}

\cl[saisine]{La saisine de la Division des Affaires Internes est \textbf{automatique} dans les cas suivants :}
\begin{itemize}
  \item tout coup de feu tiré par un agent, avec ou sans impact ;
  \item tout décès ou blessure grave survenu pendant une garde à vue ou une interpellation ;
  \item toute plainte déposée par un citoyen à l'encontre d'un agent ;
  \item tout soupçon de corruption, de fuite d'information ou de conflit d'intérêts ;
  \item tout usage de la force signalé comme disproportionné ;
  \item tout accident grave impliquant un véhicule de service ;
  \item toute altération ou disparition d'un enregistrement ou d'une pièce de procédure.
\end{itemize}

\cl{La division peut également être saisie par l'État-Major, par tout superviseur, ou par tout agent signalant un manquement dont il a été témoin.}

\sstitre{Plaintes citoyennes}

\cl[plainte]{Toute personne peut déposer une plainte à l'encontre d'un agent, au poste de police ou par formulaire écrit. Cette plainte est enregistrée et transmise sans filtrage.}

\begin{interdit}[Traitement des plaintes]
Il est interdit de dissuader une personne de déposer plainte, de refuser d'enregistrer une plainte, de lui demander de motiver sa démarche, de la différer ou de subordonner son enregistrement à une quelconque condition. Un refus d'enregistrement est en lui-même une faute lourde.
\end{interdit}

\sstitre{Procédure}

\begin{enumerate}
  \item \stress{Accusé de réception} : le plaignant ou le signalant est informé de l'enregistrement de sa démarche.
  \item \stress{Mesures conservatoires}, le cas échéant : retrait administratif, suspension du port d'arme, réaffectation hors terrain pour la durée de l'enquête. Ces mesures ne constituent pas des sanctions.
  \item \stress{Enquête} : relevé des enregistrements de caméra-piéton, audition des agents et des témoins, examen des rapports et des communications radio.
  \item \stress{Conclusions} : transmises au Chief. L'enquêteur ne communique jamais directement ses conclusions à l'agent mis en cause.
  \item \stress{Décision} : prononcée par l'autorité compétente et notifiée par écrit à l'agent, ainsi qu'au plaignant dans son principe.
\end{enumerate}

\cl[issues]{L'enquête peut donner lieu à : un classement sans suite, un rappel à la règle, une sanction disciplinaire dans les conditions du présent titre, ou la transmission du dossier à l'autorité judiciaire.}

\cl[corruption]{La corruption fait l'objet d'une tolérance nulle. La transmission d'informations de service à un civil non habilité ou à une organisation criminelle entraîne la révocation immédiate et la délivrance d'un mandat.}

\begin{rappel}[Protection de l'agent innocenté]
Lorsqu'une enquête établit qu'un agent a été mis en cause à tort, la Division des Affaires Internes en informe expressément sa hiérarchie et sa chaîne de commandement. La mention est portée à son dossier. Une accusation non fondée ne laisse aucune trace défavorable.
\end{rappel}

% ---------------------------------------------------------------------
\article{Droits de l'agent mis en cause et voies de recours}

\cl[droits-agent]{Tout agent faisant l'objet d'une enquête interne ou d'une procédure disciplinaire bénéficie des droits suivants :}
\begin{itemize}
  \item être informé, avant tout entretien, des faits précis qui lui sont reprochés ;
  \item disposer d'un délai raisonnable pour préparer ses explications ;
  \item consulter les pièces du dossier avant l'entretien disciplinaire ;
  \item être assisté d'un agent de son choix ou d'un conseil ;
  \item présenter des observations écrites versées au dossier ;
  \item être informé par écrit de la décision et de sa motivation.
\end{itemize}

\cl[recours]{Tout agent sanctionné peut former un recours devant le Command Staff et l'État-Major et solliciter l'annulation ou la réformation de la sanction. Le recours est formé par écrit et motivé.}

\cl{Le recours n'est pas suspensif, sauf décision contraire de l'autorité saisie. L'autorité statue par décision motivée et notifiée.}

\cl[recours-abus]{Aucune mesure défavorable ne peut être prise à l'encontre d'un agent au motif qu'il a exercé une voie de recours de bonne foi.}

% ---------------------------------------------------------------------
\article{Validation et modification du règlement}

\cl[validation]{Le présent règlement, ses modifications et l'abrogation de ses dispositions sont validés par l'État-Major.}

\cl{Toute modification donne lieu à l'inscription d'une nouvelle ligne au registre des versions figurant en tête du présent document, ainsi qu'à une note de service diffusée à l'ensemble des personnels.}

\cl{Il appartient à chaque agent de s'assurer de connaître le règlement et ses dernières mises à jour.}

\cl{Le présent règlement entre en vigueur le \DATEVIGUEUR, date de sa première validation par l'État-Major.}

\vspace{16pt}
\lspdornament
\vspace{6pt}

\begin{center}
{\lspdsans\fontsize{8.6}{11}\selectfont\color{lspdgrey}%
  \addfontfeature{LetterSpace=16}\MakeUppercase{Lu et approuvé le \DATEAPPROBATION\ par}}\\[5mm]
{\lspddisp\fontsize{13}{16}\selectfont\color{lspdnavy}%
  L'État-Major du Los Santos Police Department}
\end{center}

\vspace{2mm}
\begin{center}
\begin{minipage}{0.96\linewidth}\centering
\blocSignature{\signatureChief}{Arthur J. Rhodes}%
  {Chief of Police}{Badge n\textsuperscript{o} 001}
\hfill
\blocSignature{\signatureEtatMajor}{Elena R. Vasquez}%
  {Commander, État-Major}{Badge n\textsuperscript{o} 004}
\end{minipage}
\end{center}

\vspace{14pt}
\begin{note}[Note terminologique]
Dans l'ensemble du présent règlement, les termes « agent », « police officer », « officier », « employé », « personnel » et « membre » désignent indifféremment l'ensemble des employés du Los Santos Police Department.
\end{note}
